\chapter{方法选择与完整计算工作流}\label{ch:22}
\begin{chaptersummary}
从状态表示、尺度约化与闭合，到空间离散和网格，全书已形成完整的建模链。本章以目标物理量统筹这些选择，将模型、约化、空间、时间与代数误差分层检验，再把守恒、稳定和收敛组织为可执行流程。最终方法组合由目标量误差、结构保持和验证证据共同确定。
\end{chaptersummary}
\learninggoals{为给定物理目标选择模型和离散；区分五类误差；设计空间、时间和容差独立变化实验；区分数值验证与物理确认。}
\studyroute{这里组织验证方案，略去具体实验不确定度、误差估计器和项目级数据管理；代数求解与计算架构见第 0 卷。}{verification and validation; grid convergence; temporal convergence; model error; reduction error; quantity of interest}{为一个波动或输运问题定义目标量，分别改变模型、网格、次数、步长与代数容差，建立能够区分误差来源的实验矩阵。}
\section{方法选择与组合策略}
几何固定且壁面精度要求高时，贴体网格便于直接解析边界层。

对于复杂运动几何，若主要目标位于壁面剪切应力之外，浸没边界法可简化网格组织。

若局部尺度变化极大，可使用四叉树或八叉树自适应网格。

若边界运动幅度较小，可使用动网格。

若多个刚体发生大幅相对运动，可使用重叠网格。

实际工业软件常组合使用：
\[
\text{贴体边界层}
+
\text{自适应远场},
\]
\[
\text{八叉树背景网格}
+
\text{浸没边界},
\]
\[
\text{重叠网格}
+
\text{局部动网格}.
\]

\begin{askai}
AI 展开问题：
\begin{quote}
给定一个带旋转叶片和远场声波传播的计算问题，请比较贴体网格、动网格、重叠网格和自适应远场网格各自解决什么子问题。请按各子区域的需求给出组合网格策略。
\end{quote}

进一步展开：
\begin{quote}
把“几何固定性、运动幅度、边界层精度、局部尺度变化、稳定步长”做成一个方法选择表。
\end{quote}
\end{askai}
\section{完整计算工作流}
一次可靠计算通常按以下链条进行：

\begin{enumerate}
  \item 用 \(\mathrm{Kn}\)、\(\mathrm{Re}\)、\(\mathrm{Ma}\) 等尺度比选择模型；
  \item 确定守恒变量、材料模型和边界条件；
  \item 选择与几何和局部尺度匹配的网格策略；
  \item 选择满足守恒、稳定性和波动精度要求的离散格式；
  \item 根据显式或隐式推进形成代数系统；
  \item 记录代数容差与目标量对残差的敏感性，相关求解算法见第 0 卷；
  \item 分别监控模型误差、网格误差、离散误差、线性残差和非线性残差。
\end{enumerate}

残差衡量代数方程的满足程度；模型有效性与离散精度还需通过基准解、网格及时间收敛实验和物理数据验证。

\section{五类误差与验证证据}\index{验证与确认}
\begin{center}\small
\begin{tabularx}{\textwidth}{@{}l>{\raggedright\arraybackslash}X@{}}
\toprule
误差层次 & 可执行检查\\
\midrule
物理模型 & 与实验、测量或适用范围内的高保真模型比较，报告参数和数据不确定度。\\
约化与闭合 & 改变小参数，比较完整与约化模型；检验未解矩和尺度项的量级。\\
空间离散 & 固定充分小的步长与容差，分别细化 $h$、增加 $p$，比较目标量和预期阶数。\\
时间离散 & 固定充分解析的空间，减半 $\Delta t$，检查时间阶数与相位累积。\\
代数求解 & 收紧残差容差，观察目标量变化是否已小于空间和时间误差。\\
\bottomrule
\end{tabularx}
\end{center}
几何近似、速度域截断、随机抽样和边界模型还应在相应层次单列。
代数残差需通过算子稳定性或目标量敏感性与解误差联系。

\section{守恒、稳定与收敛的实验}
在周期或封闭无源系统中，记录总质量、动量与能量；
有物理源和边界时，记录输入输出及耗散的平衡缺口。
能量估计或离散熵给出允许的演化范围，波动问题还需测量相位与群速度。
电磁和 MHD 问题同时记录离散散度。

光滑解析解可检验预期阶数，制造解可指定源项和边界来覆盖全部算子。
间断、壁面层和多尺度问题需采用适合其正则性和目标量的误差指标。
数值验证考察给定模型的求解误差；物理确认考察模型对目标现象的描述能力。

\conceptbridge{残差、守恒和网格收敛都通过后，模型还需什么证据？}
{这些检查针对特定离散和模型的数学表现。物理确认还需目标条件下的测量或适当基准，
并评估本构、边界和参数的不确定度。各层证据共同限定计算结论的可信范围。}
{数值验证 $\longrightarrow$ 模型与数据比较 $\longrightarrow$ 目标量的可信范围。}

\aiexercise{为周期 Alfv\'en 波或低 Mach 剪切波建立误差表，分别列出网格、次数、步长、容差、守恒缺口和目标量误差。设计能够识别时间误差掩盖空间收敛的对照组。}
{将目标换成壁面热流或湍流输运，列出新增的碰撞、本构、边界和统计假设，以及可以独立检验这些假设的数据。}

\section{全书的依赖链}
\[
\text{Liouville 与统计表示}
\longrightarrow\text{动理学与连续模型}
\longrightarrow\text{尺度约化与闭合}
\]
\[
\longrightarrow\text{函数空间与弱形式}
\longrightarrow\text{残差、通量与高阶离散}
\longrightarrow\text{网格与验证}.
\]
保留的状态决定所需闭合，连续算子与目标结构决定适合的离散约束，
目标量与验证证据再决定计算分辨率和方法组合。
